\section{Reasoning：从“会想”到“会做”}

\subsection{Reasoning 的工程目标}
Yu Su 在课程中反复强调：Reasoning 的最终目标不是展示推理链，而是提高行动质量（Action Quality）。如果推理不能带来更好的动作序列、更低错误率或更高样本效率，那么该推理就是无效开销。

\begin{figure}[H]
\centering
\includegraphics[width=0.84\textwidth]{slide-02-reasoning-action.jpg}
\caption{视频约 00:21:50 讲义页：Reasoning 与 Acting 的耦合关系是课程主轴}
\end{figure}

\begin{importantbox}{Reasoning for Acting}
\begin{itemize}
    \item 用推理提升动作选择质量，而不是增加表面思考长度；
    \item 推理应当可触发策略修正（replan）和工具切换（tool switch）；
    \item 推理的价值最终体现在成功率、成本与鲁棒性三项指标上。
\end{itemize}
\end{importantbox}

\subsection{从 CoT 到可执行推理}
课程把 CoT、ReAct、Self-Reflection、Meta-Reasoning 放在同一连续谱中：它们本质上都是 ``在动作前后注入结构化思考''。差异不在于是否推理，而在于推理如何连接环境状态和动作接口。

\begin{knowledgebox}{四类常见推理动作}
\begin{enumerate}
    \item 任务分解：将复杂目标拆成可验证子目标；
    \item 证据对齐：用外部检索/工具输出修正内部假设；
    \item 失败归因：识别错误来源（知识缺失、工具失败、规划错误）；
    \item 策略改写：基于归因结果重写后续动作计划。
\end{enumerate}
\end{knowledgebox}

\begin{warningbox}{推理 Token 并非越多越好}
在没有有效验证器与边界约束时，长推理链可能导致：
\begin{itemize}
    \item 成本爆炸；
    \item 置信幻觉放大；
    \item 错误路径被 ``解释得更像对''。
\end{itemize}
\end{warningbox}

\subsection{本章小结}
Reasoning 是 Agent 的决策增益机制。真正重要的不是 ``有没有推理''，而是 ``推理是否改变了更优行动''。

